\section{The Adios Design and Its Reported Results}
\label{sec:adios}
To establish a clear baseline, we first outline the architecture and empirical claims of the EuroSys '25 Adios paper \cite{adios} that inspired our simulation study.

\subsection{Motivation and Baseline Limitations}
Yoon et al.\ benchmarked DiLOS, the leading state-of-the-art busy-waiting memory disaggregation system, on an 8-core compute node using a random array-access microbenchmark with 20\% local DRAM caching. Their measurements revealed an early throughput ceiling: request throughput flattened out at approximately 1.4~MRPS, while tail latency (P99) surged starting at 1.3~MRPS. Inspecting per-request latency breakdowns showed that queueing delays in front of spinning cores caused this tail blowout, while the 100~GbE RDMA network sat barely half-utilized ($\approx 50\%$). Furthermore, adding Concord-style 5~$\mu$s CPU preemption failed to help, confirming that spinning threads cannot be effectively sliced because the delay is network transit rather than active instruction execution.

\subsection{Architectural Components of Adios}
Adios eliminates spinning through several integrated design choices:
\begin{itemize}
\item \textbf{OSv Unikernel Base}: Application code, memory paging, and the task scheduler co-exist within one ring-0 address space, turning scheduler invocations into local function calls without hardware privilege-level crossings.
\item \textbf{Ultra-Lightweight Unithreads}: Rather than heavy POSIX pthreads, Adios wraps each incoming request into a minimal thread context requiring only six saved registers (80 bytes total). A shared universal stack collapses kernel and user stacks, capping memory per request at 4~KiB and slashing switch cost to 40 CPU cycles (compared to 191 cycles in Shinjuku \cite{shinjuku}).
\item \textbf{Yield-on-Fault Paging}: On a page fault, the thread posts an RDMA read and yields the core. The worker processes other waiting requests, polling completion queues between tasks. When the page lands, the faulting thread resumes, maps the memory, and completes.
\item \textbf{Proactive Reclaiming}: A dedicated core runs an eviction daemon that keeps free page frames continuously available, preventing faults from blocking on synchronous evictions.
\item \textbf{Page-Fault-Aware Dispatching}: Because yielding workers can accumulate different numbers of active RDMA fetches, per-worker queue pairs can develop unbalanced queues. As shown in Algorithm~\ref{alg:dispatch}, the central dispatcher assigns incoming requests to whichever idle worker currently has the fewest in-flight fetches.
\item \textbf{Asynchronous Reply Delegation}: Sending response packets back to clients is offloaded to the dispatcher's completion queue, avoiding worker spinning during transmission.
\end{itemize}

\begin{algorithm}[t]
\caption{Page-fault-aware dispatching (our rendering of the policy in \cite{adios}).}
\label{alg:dispatch}
\begin{algorithmic}[1]
\small
\Procedure{Dispatch}{$Q$, $W$}
\State $I \gets \{w \in W : w \text{ is idle}\}$
\State sort $I$ by outstanding fetches of $w$, ascending
\While{$Q \neq \emptyset$ \textbf{and} $I \neq \emptyset$}
  \State $r \gets \text{pop}(Q)$;\; $w \gets \text{pop}(I)$
  \State send $r$ to $w$
\EndWhile
\EndProcedure
\end{algorithmic}
\end{algorithm}

\subsection{Published Evaluation Results}
Using 8 worker cores, 1 dispatcher, and 1 reclaimer across an InfiniBand/RoCE 100~GbE network, Adios demonstrated impressive real-world improvements over DiLOS (summarized in Table~\ref{tab:paper}). Two specific findings motivated our simulation: First, when local memory drops from 100\% down to 10\%, Adios's throughput degrades by only 25\%, whereas DiLOS loses 60\%. Second, under very light load (0.2~MRPS), Adios exhibits a minor latency regression ($\approx 6\%$ higher P99.9), because paying context-switch penalties offers zero concurrency benefit when no other requests are queued.

\begin{table}[t]
\caption{Selected results reported for Adios \cite{adios}}
\label{tab:paper}
\centering
\footnotesize
\begin{tabular}{p{3.9cm}p{3.2cm}}
\toprule
Quantity & Reported value \\
\midrule
DiLOS throughput plateau (microbenchmark) & about 1.4 MRPS \\
Adios throughput growth (microbenchmark) & up to about 2.5 MRPS, $\approx 1.5\times$ DiLOS \\
RDMA link utilization & about 50\% (DiLOS) vs.\ up to 82\% (Adios) \\
Queueing delay at 1.3 MRPS & $16.3\times$ (P99) and $36.8\times$ (P99.9) smaller \\
Context switch & 40 vs.\ 191 cycles (Shinjuku) \\
PF-aware dispatch gain on P99.9 & up to 7.5\% (Memcached), 27\% (RocksDB) \\
Real applications vs.\ DiLOS & throughput $1.07$--$1.64\times$; P99.9 $1.99$--$10.89\times$ lower \\
\bottomrule
\end{tabular}
\end{table}

\subsection{Open Questions and Simulation Scope}
The Adios authors openly acknowledge several practical boundaries. Single-dispatcher queueing does not scale beyond roughly 10--12 workers due to central lock contention. Two whole CPU cores must be sacrificed exclusively for dispatching and page reclaiming. Most importantly, yielding benefits rely on having a pool of concurrent requests ready to fill the idle fetch intervals; single-threaded or CPU-bound applications see negligible speedups. Table~\ref{tab:modelled} clarifies which of these architectural features our discrete-event simulator directly captures versus abstracts away.

\begin{table}[t]
\caption{Which Adios mechanisms the simulator represents}
\label{tab:modelled}
\centering
\footnotesize
\begin{tabular}{p{2.5cm}p{1.1cm}p{3.5cm}}
\toprule
Mechanism & Modelled? & How \\
\midrule
Yield on fault & Yes & Fetch issued, worker serves other requests \\
Context-switch cost & Yes & Parameter $t_s$, swept from 0.02 to 4~$\mu$s \\
Single central queue & Yes & One FIFO for both policies \\
PF-aware dispatch & Yes & Fewest outstanding fetches first \\
RDMA queue pairs & Partly & One FIFO server per worker, plus a shared link \\
Local cache & Partly & Fixed hit rate; optional LRU with Zipf access \\
Unithread layout, universal stack & No & Only the switch cost is kept \\
Polling delegation & No & Reply path is not simulated \\
Reclaimer thread & No & Free frames are always available \\
Real NIC and kernel behaviour & No & Fixed service times \\
\bottomrule
\end{tabular}
\end{table}
